\documentclass[10pt]{article}

\usepackage[utf8]{inputenc}

\usepackage[T1]{fontenc}

\usepackage{graphicx}

\usepackage[export]{adjustbox}

\graphicspath{ {./images/} }

\usepackage{amsmath}

\usepackage{amsfonts}

\usepackage{amssymb}

\usepackage[version=4]{mhchem}

\usepackage{stmaryrd}

\usepackage{CJKutf8}



\DeclareUnicodeCharacter{0131}{$\imath$}



\begin{document}

\includegraphics[max width=\textwidth, center]{2023_07_09_d44b24b18c615cb15de8g-1(3)}

ф 0 н д A. O., "C. r. Acad. sci.», 1929, t. 189, p. 1224-28; [7] S i e g' e l C., "Abhandl. Preuss.' Akad. Wiss., Phys. Cl.", 1929, № 1 , S. $1-70 ;[8] \Gamma$ е ль ф о н д А. O., "Докл. АН CCCP》, 1934, T. 2, c. $1-6$; [9] S ch n e i d e r T h., "J. reine und angew. Math.», 1934, Bd 172, S. 65-74; [10] B a k e r A., Transcenden-



\includegraphics[max width=\textwidth, center]{2023_07_09_d44b24b18c615cb15de8g-1(4)}

[12] $\Phi$ е ль д м а и Н. И., Седьмая проблема Гильберта, М., 1982. В. Г. Спринджук.



ТРАНСЦЕНДЕНТНОСТИ МЕРА - функция, характеризугщая отклонение заданного трансцендентного числа от множества алгебраич. чисел ограниченной степени и ограниченной высоты ггри изменении границ для әтих параметров. Для трансцендентного $\omega$ и натуральных $n$ и $H$ Т. м. есть



\[

w_{n}(\omega ; H)=\min |P(\omega)|,

\]



где минимум беротся по всем ненулевым целочисленным многочленам $P(x)$ степени не более $n$ и высоты не более Н. Из Дирихле припцииа «ящиков» следует, что всегда



\[

w_{n}(\omega ; I)<c_{1}^{n} H^{-n},

\]



где $c_{1}$ зависит только от $\omega$. Во многих случаях вместе с доказательством трансцендентности числа $\omega$ удается получить оценку снизу для Т. м. в терминах степенной, логарифмическої и әкспонешцальной функций от $n$ и $H$. Напр., метод Эрмита доказательства трансцендентности $e$ позволяет получить неравенство



\[

w_{n}(e ; H)>H^{-n-\frac{c_{2} n^{2} \ln n}{\ln \ln H}} \text {, }

\]



где $c_{2}>0$ - абсолютная постоянная и $H \geqslant H_{0}(n)$. При любых фиксированных $n$ и $\varepsilon>0$



\[

w_{n}(\omega ; H)>c_{3} H^{-n-\varepsilon}, c_{3}=c_{3}(\omega ; n, \varepsilon)

\]



для почти всех (в смысле Јебега) действительных чисел $\omega$ (см. Малера проблема). Трансцендентные числа могут быть классифицированы на основе различий в асимптотич. поведении $w_{n}(\omega ; H)$ при неограниченном изменении $n$ и $H$ (см. [3]).



Лит.: [1] Г е ль ь о і д А. О., Трансцендентные и алгебраические тисла, М., 1952; [2] C i j s o u w P. L., Transcendence measures, Amst., 1972 (Diss.); [3] B a k e r A.,' Transcendenta1 number theory, L., 1975. B. Г. Спринджук.



ТРАПЕЦЙИ ФОРМУЛА - частный случай $H_{\text {ьюто- }}$ на- Котеса квадратурной форлуль, в к-рої берется два узла:



\[

\int_{a}^{b} f(x) d x \cong \frac{b-a}{2}[f(a)+f(b)]

\]



Если подинтегральная фуншция $f(x)$ сильно отличается от линейной, то формула (1) дает малую точность. Промежуток $[a, b]$ разбивается на $n$ частичных промежутков $\left[x_{i}, x_{i+1}\right], \quad i=0,1, \ldots, n-1$, длины $h=$ $=(b-a) / n$ и для вычисления интеграла по промежутку $\left[x_{i}, x_{i+1}\right]$ используется Т. к. ф.



\[

\int_{x_{i}}^{x_{i+1}} f(x) d x \cong \frac{h}{2}\left[f\left(x_{i}\right)+f\left(x_{i+1}\right)\right] .

\]



Суммирование левой п правой части әтого приближенного равенства по $i$ от 0 до $n-1$ приводит к составної T. к. Ф.:



\[

\begin{gathered}

\int_{a}^{b} f(x) d x \cong h\left[\frac{f(a)}{2}+f\left(x_{1}\right)+\right. \\

\left.+f\left(x_{2}\right)+\ldots+f\left(x_{n-1}\right)+\frac{f(b)}{2}\right],

\end{gathered}

\]



где $x_{j}=a+j h, j=0,1,2, \ldots, n$. Квадратурную формулу (2) также называют Т. к. ф. (без добавления слова составная). Алгебраич. степень точности квадратурной формулы (2), как и формулы (1), равна 1. Квадратурная формула (2) точна для тригонометрич. функций



$\cos \frac{2 \pi}{b-a} k x, \sin \frac{2 \pi}{b-a} k x, k=0,1,2, \ldots, n-1$. В случае $b-a=2 \pi$ формула (2) точна для всех тригонометрич. полиномов порядка не выше $n-1$, более того, ее тригонометрич. стенень точности равна $n-1$.



Если подинтегральная функция $f(x)$ дважды непрерывно диффференцируема на $[a, b]$, то погрешность $R(f)$ квадратурной формулы (2) - разность между интегралом и квадратурної суммой - имеет представление



\[

R(f)=-\frac{b-a}{12} h^{2} f^{\prime \prime}(\xi),

\]



где $\xi$ - нек-рая тотка промежкутка $[a, b]$.



\begin{enumerate}

  \setcounter{enumi}{10}

  \item П. Мысовских. ТРАПЕЦИЯ - выпуклый четырехугольник, у к-рого две стороны параллельны, а две другие - непараллельны (см. рис.). Параллельные стороны наз. о с н ов а н и я м и Т., а непараллельные - ее б о к о в ы м и с т ор о н а ми; отрезок, соединягощий середины боковых сторол Т.,- еe c p e л н е ї ли н и ё̈̈: средняя линия параллельна ос-

\end{enumerate}



\includegraphics[max width=\textwidth, center]{2023_07_09_d44b24b18c615cb15de8g-1(1)}

нованиям $\mathrm{T}$. и равна их полусумме. Площадь $\mathrm{T}$. равна произведению средней лишии на высоту Т. или половине произведения диагоналей па синус угла между ними. Т., боковые стороны к-роӥ равıы между собой, наз. р а в н о б о ч н о \begin{CJK}{UTF8}{mj}立\end{CJK}.



TPETH T PAEBAG задач для лифференциальных уравнений с тастными производными. Пусть, напр., в ограниченной области $\Omega$, в каждой точке границы $\Gamma$ к-рой существует нормаль, задано эллиптич. уравнение 2-го порядка



\[

\begin{gathered}

L u=\sum_{i, j=1}^{n} a_{i j}(x) \frac{\partial^{2} u(x)}{\partial x_{i} \partial x_{j}}+\sum_{i=1}^{n} b_{i}(x) \frac{\partial u(x)}{\partial x_{i}}+ \\

+c(x) u(x)=f(x),

\end{gathered}

\]



где $x=\left(x_{1}, x_{2}, \ldots, x_{n}\right), n \geqslant 2$. Т. к. з. для уравнения (ж) в области $\Omega$ наз. следуюццая задача: из множества всех решений и $u(x)$ уравнения (:) выделить те, к-рые в каждой граничной точке имеют производные по внутренней конормали $N$ и удовлетворяют условию



\[

\frac{\partial u(x)}{\partial N}+\alpha(x) u(x)=v(x), x \in \Gamma,

\]



где $\alpha>0$ и $v-$ заданные непрерывные на $\Gamma$ функции.



A. В. ІІвангов.



ТРЕУГОЛЬНАЯ ГРУППА, Т $\mathrm{p}$ И а н Г У л и $\mathrm{p}$ У ема я гр у п п а, - то же, что Ли вполне разрецимая әр ynna.



ТРЕУ ГОЛЬНАЯ МАТРИЦА - квадратная матрица, у к-роӥ все элементы, расположеннце ниже (или выше) главной диагонали, равны нулю. В первом случае матрица наз. в е р х н е й т р е у г о л ь н о й м ат р и ц ї, во втором- - и й п е й т р е уг о льн о й м а т р и ц е й. Определитель Т. м. равен произведениі всех ее диагональных әлементов.



ТРЕУТО О. А. Иванова.



\includegraphics[max width=\textwidth, center]{2023_07_09_d44b24b18c615cb15de8g-1(2)}

с т и - три точки (вершины) и три отрезка прямых (стороны) с концами в этих точках. Иногда при определении Т. к нему относят и выпуклуіо часть глоскости, к-рая огранитена сторонами $\mathrm{T}$.



Понятие Т. вводится и в мпогообразиях, отличных от евклидовой плоскости. Т. обычно определяется как три точки и три отрезка геодезических с концами в этих точках. Таковы, напр., сферич. треугольниг в сферической геометрии, треугольник в плоскости Лобачевского (см. Неевклидовь геометрии).



\includegraphics[max width=\textwidth, center]{2023_07_09_d44b24b18c615cb15de8g-1}

Новые встречи с геометрией, пер. с англ., М., 1978; [2] К о к ст е $\mathrm{p}$ Х. С. М., Введепие в геометрию, пер. с англ., М., 1966; [3] Зете ль С. И., Ноовая геометрия треугольника, 2 изд., М., 1962; [4] А д а м а р ЖК., Элементарная геометрия, ч. 1. Планиметрия, пер. с франц., ч изд., М., 1957; [5] Е Ф р ем о в Д., Новая геометрия треугольника, одесса, 1902 .





\end{document}